\chapter{Sprite caching and partial paint}
\label{ch:sprite}

\begin{aside}
A box you draw a hundred times a second is a box you should
draw once and stamp. Sprite caching is the renderer's optional
``don't repaint that thing'' optimisation.
\end{aside}

\section{The pattern}

A complex widget (gradient title bar, decorated card, ASCII
logo) paints the same cells every frame. The cells are computed
from props the renderer has seen before.

A sprite cache memoises: key on the props, store the cell
array, recall on hit.

\section{Cache structure}

\begin{lstlisting}
struct Sprite {
    int w, h;
    std::vector<Cell> cells;
};
std::unordered_map<Key, Sprite> cache_;
\end{lstlisting}

The key is a hash of the widget's render-relevant inputs:
text, style, width, height. On hit, blit the cached cells
into the back buffer.

\section{When to cache}

Caching pays off when:

\begin{itemize}[leftmargin=*]
  \item Paint cost is high (computed gradients, glyph
    composition).
  \item The widget renders repeatedly with the same props.
  \item Memory budget allows.
\end{itemize}

Caching hurts when:

\begin{itemize}[leftmargin=*]
  \item Props change every frame.
  \item Widget is cheap to paint anyway.
  \item Cache eviction thrashes.
\end{itemize}

\section{Eviction}

LRU is the standard policy. Size the cache to a memory
budget: number of cells, not number of entries.

\maya{}'s default budget is 256 KB. A modest UI uses
~10--20 KB; large media-heavy UIs may need more.

\section{Invalidation}

When a prop changes, the cache key changes; the old entry
sits in cache until evicted by LRU. No explicit
invalidation needed.

If a global setting changes (theme, font, locale) that
should invalidate everything, call \code{cache.clear()}
explicitly.

\section{Compositing onto the back buffer}

\begin{lstlisting}
void blit(Painter& p, Rect dst, const Sprite& s) {
    for (int dy = 0; dy < s.h; ++dy)
        for (int dx = 0; dx < s.w; ++dx)
            p.put(dst.x + dx, dst.y + dy,
                  s.cells[dy * s.w + dx]);
}
\end{lstlisting}

A nested loop; cheap. The cache replaces the widget's own
paint method.

\section{Partial paint}

A stronger form: only paint what changed since last frame.
Requires per-widget dirty tracking and a damage rectangle
list to feed the diff.

\maya{} does not do partial paint. The trade-off:
bookkeeping cost vs cell-diff cost. The cell diff already
finds unchanged regions cheaply; partial paint would
duplicate the work.

For very large screens (>200x100), partial paint becomes
attractive. \maya{}'s sweet spot is normal terminal sizes.

\section{Pre-rendered assets}

For static elements (logo, splash screen, decorative ASCII
art), compute the cells once at startup, store in static
storage, blit each frame. This is sprite caching's degenerate
case --- key is the constant ``true''.

\maya{} provides \code{static\_sprite(...)} as a helper.

\section{Recap}

\begin{recap}
\begin{itemize}[leftmargin=*]
  \item Sprite cache memoises expensive paint.
  \item LRU eviction with a memory budget.
  \item No explicit invalidation; key change is sufficient.
  \item Partial paint not implemented; cell diff suffices
    for typical sizes.
  \item Static sprites for unchanging assets.
\end{itemize}
\end{recap}

\section*{Workshop}

\begin{enumerate}[leftmargin=*]
  \item Identify your most expensive widget. Add a sprite
    cache; measure the speedup.
  \item Stress test eviction by changing a widget's prop
    every frame; watch cache thrash.
  \item Pre-render an ASCII logo as a static sprite.
\end{enumerate}
